\chapter{形变理论}\label{src:C}
\setlength{\emergencystretch}{4em}
\hypersetup{pageanchor=false}
\AddToHook{shipout/foreground}[appendix-c-index-pages]{\put(0,0){\makebox[0pt][l]{\hypertarget{page.\thepage}{}}}}
\begin{epigraph}
形变理论是一个很难写清楚的主题. 现有文献可以分为两类.

有些书考虑一个非常特殊的情形, 并以不难理解的方式详细展开. 我想到的最佳例子是关于奇点形变的 \cite{Art76}, 关于流形形变理论的 \cite{Kod86}, 以及关于代数曲面上曲线的 Hilbert 概形的 \cite{Mum66a}. 这些作者暗含的主张是, 如果真正理解了这些特殊情形, 处理更一般的情形就不太困难. 我的经验表明这确实如此, 但需要付出相当多的艰苦努力.

另一种方法是在最一般的背景下发展形变理论, 并证明那些一劳永逸地解决这一主题的定理. 这里最好的例子是 \cite{Ill71,Bin87}. 遗憾的是, 一般形变理论似乎需要一整套令人望而生畏的工具, 这使这些论著很难阅读.
\epigraphauthor{János Kollár \cite[第 $6$ 页]{Kol96}}
\end{epigraph}
形变理论研究模空间 $\mathcal M$ 在对象 $E_0\in\mathcal M(\Bbbk)$ 附近的局部几何. 我们主要关注以下三个形变问题:
\begin{enumerate}[(A)]
\item 域 $\Bbbk$ 上固定射影概形 $X$ 中闭子概形 $Z_0$ 的嵌入形变. 此时模问题为 Hilbert 函子 $\Hilb^P(X)$, 对象为 $E_0=[Z_0\subseteq X]\in\Hilb^P(X)(\Bbbk)$.
\item 域 $\Bbbk$ 上概形 $E_0$ 的形变. 我们的出发例子是 $E_0$ 为光滑曲线的情形, 此时模问题为 $\mathcal M_g$, 且 $[E_0]\in\mathcal M_g(\Bbbk)$; 更一般地, 若 $E_0$ 为固有曲线, 则模问题为所有曲线组成的叠 $\mathcal M_g^{\mathrm{all}}$.
\item 域 $\Bbbk$ 上固定射影概形 $X$ 上凝聚层 $E_0$ 的形变. 模问题为 $\Coh(X)$, 且 $[E_0]\in\Coh(X)(\Bbbk)$.
\end{enumerate}
形变理论提供了由局部到整体考察模的视角. 放大 $E_0\in\mathcal M(\Bbbk)$ 附近的区域, 我们依次研究 $\mathcal M$ 在 $E_0$ 处的一阶邻域, $E_0$ 的高阶形变, $E_0$ 的形式邻域, 最后是 $E_0$ 的平展或光滑邻域. 开始之前, 先简要介绍本附录的七节内容.
\begin{enumerate}[(1)]
\item $E_0$ 的一阶形变是双数环 $\Bbbk[\epsilon]:=\Bbbk[\epsilon]/(\epsilon^2)$ 上的对象 $E\in\mathcal M(\Bbbk[\epsilon])$, 连同同构 $\alpha:E_0\to E|_{\Spec\Bbbk}$; 换言之, 是交换图
\[
\begin{tikzcd}
\Spec\Bbbk\arrow[r,hook]\arrow[dr,"{[E_0]}"']&\Spec\Bbbk[\epsilon]\arrow[d,"{[E]}"']\\
&\mathcal M,
\end{tikzcd}
\]
它使我们能够将 $E$ 看成 $\mathcal M$ 在 $E_0$ 处的切向量. 我们在 \S\bookref{C.1} 中分类问题 (A)--(C) 的一阶形变.
\item 给定剩余域为 $\Bbbk$ 的 Artin 局部 $\Bbbk$-代数满射 $A'\twoheadrightarrow A$, 以及对象 $E\in\mathcal M(A)$ 和同构 $E_0\to E|_{\Spec\Bbbk}$, $E$ 在 $A'$ 上的形变是对象 $E'\in\mathcal M(A')$ 连同同构 $\alpha:E\to E'|_{\Spec A}$. 图示为交换图
\[
\begin{tikzcd}
\Spec\Bbbk\arrow[r,hook]\arrow[drr,"{[E_0]}"',bend right=12]&\Spec A\arrow[r,hook]\arrow[dr,"{[E]}"']&\Spec A'\arrow[d,"{[E']}"']\\
&&\mathcal M.
\end{tikzcd}
\]
在 \S\bookref{C.2} 中, 我们对问题 (A)--(C) 确定 $E$ 在 $A'$ 上的形变 $E'$ 存在的障碍, 并在障碍消失时分类所有这样的形变.
\end{enumerate}
前两节对于研究 $\overline{\mathcal M}_g$ 和 $\cBun(C)$ 至关重要, 例如建立光滑性并计算维数; 见定理 \bookref{6.4.16}. 由于形变理论的研究与叠和模不可分割, 为求完整, 我们还收入了 \S\bookref{C.3}--\bookref{C.7}.
\begin{enumerate}[(1)]\setcounter{enumi}{2}
\item 在 \S\bookref{C.3} 中, 我们概述余切复形, 特别说明它如何控制无穷小形变理论, 从而重新得到 \S\bookref{C.1}--\bookref{C.2} 的一些结果.
\item 给定完备 Noether 局部 $\Bbbk$-代数 $(R,\mathfrak m)$, $E_0$ 在 $R$ 上的形式形变是一族相容的形变 $E_n\in\mathcal M(R/\mathfrak m^{n+1})$. 若 Artin 环加厚上的每个形变都经由某个 $E_n$ 分解, 则称形式形变 $\{E_n\}$ 为\textit{万有的} (精确定义见定义 \bookref{C.4.2}). Rim--Schlessinger 判别法 (定理 \bookref{C.4.6}) 给出 $E_0$ 的万有形变 $\{E_n\}$ 存在的判据; 在 \S\bookref{C.4} 中, 我们对问题 (A)--(C) 验证这些判据.
\item 若存在延拓 $\{E_n\}$ 的对象 $\widehat E\in\mathcal M(R)$, 则称 $(R,\mathfrak m)$ 上的形式形变 $\{E_n\}$ 为\textit{有效的}; 换言之, 存在交换图
\[
\begin{tikzcd}[column sep=small]
\Spec R/\mathfrak m\arrow[r,hook]\arrow[drrrr,"{[E_0]}"',bend right=12]&\Spec R/\mathfrak m^2\arrow[r,hook]\arrow[drrr,"{[E_1]}"',bend right=8]&\Spec R/\mathfrak m^3\arrow[r,hook]\arrow[drr,"{[E_2]}"']&\cdots\arrow[r,hook]&\Spec R\arrow[d,"{[\widehat E]}"']\\
&&&&\mathcal M.
\end{tikzcd}
\]
在 \S\bookref{C.5} 中, 我们证明 Grothendieck 存在定理 (\bookref{C.5.3}), 并说明它如何推出问题 (A)--(C) 的形式形变有效.
\item 给定 $R$ 上的有效万有形式形变 $\widehat E$, Artin 代数化 (\bookref{C.6.8}) 保证存在有限型 $\Bbbk$-概形 $U$, 点 $u\in U(\Bbbk)$ 及对象 $E\in\mathcal M(U)$, 使得 $R\simeq\widehat{\mathcal O}_{U,u}$, 且对所有 $n$ 有 $\widehat E|_{\Spec R/\mathfrak m^{n+1}}\simeq E|_{\Spec R/\mathfrak m^{n+1}}$.
\item Artin 代数性公理 (\bookref{C.7.1} 和 \bookref{C.7.4}) 提供验证模问题 $\mathcal M$ 的代数性的判据. 具体而言, 它给出保证上述态射 $[E]:U\to\mathcal M$ 在 $E_0$ 的某个开邻域上光滑的条件.
\end{enumerate}
关于形变理论的其他代数论述, 推荐 \cite{Art76}, \cite[\S I.$2$]{Kol96}, \cite{Vis97}, \cite[\S $6$]{FGI+05}, \cite{Ser06}, \cite{Nit09}, \cite{Har10} 及 \cite[Tag 0ELW]{SP}. 关于流形的形变, 推荐 \cite{Kod86}; 关于余切复形的详尽论述, 推荐 \cite{Ill71,Ill72}.
